\documentclass[11pt,a4paper]{article}

% ------------------------
% Encoding & fonts
% ------------------------
\usepackage[T1]{fontenc}
\usepackage[utf8]{inputenc}
\usepackage{lmodern}

% ------------------------
% Page layout
% ------------------------
\usepackage[a4paper,margin=1in]{geometry}
\usepackage{setspace}
\onehalfspacing

% ------------------------
% Math packages
% ------------------------
\usepackage{amsmath,amssymb,bm,mathtools}
\usepackage{siunitx}

% ------------------------
% Graphics & tables
% ------------------------
\usepackage{graphicx}
\usepackage{subcaption}
\usepackage{booktabs}
\usepackage{multirow}

% ------------------------
% Algorithms
% ------------------------
\usepackage{algorithm}
\usepackage{algpseudocode}

% ------------------------
% Colors & highlighting
% ------------------------
\usepackage{xcolor}
\definecolor{linkblue}{RGB}{0,90,160}

% ------------------------
% References & citations
% ------------------------
\usepackage[numbers,sort&compress]{natbib}

% ------------------------
% Glossary & acronyms
% ------------------------
\usepackage[acronym]{glossaries}
\makeglossaries

\newacronym{cnn}{CNN}{Convolutional Neural Network}
\newacronym{snr}{SNR}{Signal-to-Noise Ratio}
\newacronym{hsi}{HSI}{Hyperspectral Imaging}

\newglossaryentry{eventcamera}{
  name={event camera},
  description={A vision sensor that outputs asynchronous brightness change events}
}

% ------------------------
% Hyperlinks & clever refs
% ------------------------
\usepackage[
  colorlinks=true,
  linkcolor=linkblue,
  citecolor=linkblue,
  urlcolor=linkblue
]{hyperref}
\usepackage{cleveref}

% ------------------------
% Title information
% ------------------------
\title{\textbf{A Sample Paper for Testing the PaperAgent LaTeX Environment}}
\author{Rongzhou Chen}
\date{\today}

% ========================
% Document
% ========================
\begin{document}

\maketitle

\begin{abstract}
This document serves as a compact yet comprehensive \LaTeX{} testbed for the
PaperAgent environment. It demonstrates commonly used packages for mathematical
typesetting, figures, algorithms, references, and glossary management.
\end{abstract}

\section{Introduction}

Modern computational imaging and machine learning research often relies on
precise mathematical notation, reproducible algorithms, and structured
cross-referencing. Tools such as \gls{cnn} and \gls{hsi} are now ubiquitous in
the field of imaging systems.

An emerging sensor modality is the \gls{eventcamera}, which offers high temporal
resolution and low latency.

\section{Mathematical Formulation}

Let $\bm{x} \in \mathbb{R}^n$ denote an input signal and $\bm{W}$ a linear operator.
A common objective function is

\begin{equation}
\label{eq:loss}
\mathcal{L}(\bm{W}) =
\frac{1}{2} \lVert \bm{W}\bm{x} - \bm{y} \rVert_2^2
+ \lambda \lVert \bm{W} \rVert_1 ,
\end{equation}

where $\lambda > 0$ controls sparsity.

We report signal quality using the \gls{snr}, defined as

\begin{equation}
\mathrm{SNR} = 10 \log_{10} \left(
\frac{\lVert \bm{s} \rVert_2^2}{\lVert \bm{n} \rVert_2^2}
\right) \, \si{\decibel}.
\end{equation}

\section{Algorithm}

\begin{algorithm}[h]
\caption{Gradient Descent Optimization}
\label{alg:gd}
\begin{algorithmic}[1]
\Require Initial parameter $\bm{W}_0$, step size $\eta$
\For{$k = 0$ to $K-1$}
  \State $\bm{g}_k \gets \nabla \mathcal{L}(\bm{W}_k)$
  \State $\bm{W}_{k+1} \gets \bm{W}_k - \eta \bm{g}_k$
\EndFor
\State \Return $\bm{W}_K$
\end{algorithmic}
\end{algorithm}

\section{Experimental Results}

\begin{table}[h]
\centering
\caption{Quantitative comparison of reconstruction quality}
\label{tab:results}
\begin{tabular}{lcc}
\toprule
Method & SNR (dB) $\uparrow$ & Runtime (ms) $\downarrow$ \\
\midrule
Baseline & 28.4 & 12.1 \\
Proposed & \textbf{32.7} & 14.3 \\
\bottomrule
\end{tabular}
\end{table}

\section{Discussion}

As shown in \cref{eq:loss,alg:gd,tab:results}, the proposed formulation improves
reconstruction fidelity at a modest computational cost.

\section{Conclusion}

This sample paper validates the PaperAgent \LaTeX{} toolchain and provides a
ready-to-extend starting point for full research manuscripts.

\printglossaries

\bibliographystyle{plainnat}
\begin{thebibliography}{9}

\bibitem{goodfellow2016deep}
I.~Goodfellow, Y.~Bengio, and A.~Courville.
\newblock \emph{Deep Learning}.
\newblock MIT Press, 2016.

\bibitem{gallego2022event}
G.~Gallego et~al.
\newblock Event-based vision: A survey.
\newblock \emph{IEEE TPAMI}, 44(1):154--180, 2022.

\end{thebibliography}

\end{document}
